Terminology Locked: Physical Monetary Value Exchange

The reciprocal mirror is now strictly defined as an accounting conservation rule that governs the physical monetary value exchange between layers. All thermodynamic or “energy transfer” language is removed.

Locked Definition

\[
\mu(t) = \frac{1}{\lambda(t)}
\]

This operator records the conversion of liquid surface resources into illiquid, long-duration claims when visible growth \(\lambda(t)\) compresses.

As surface growth declines toward the ~1.5 % floor, liquid cash generation weakens.  
The conservation rule requires that the residual monetary value be absorbed into the hidden accumulator as non-liquid, long-duration claims.  
The result is a measurable shift: liquid corporate and sovereign cash flows are extracted from the surface manifold and locked into fixed infrastructure, duration-heavy assets, and related claims.

Micro Illustration (Alphabet Q2 2026)

The provided corporate data illustrate the same exchange at firm level:

Surface revenue remains strong.  
Capital expenditure rises sharply to fund long-duration technical infrastructure.  
Free cash flow compresses as liquid resources are converted into fixed assets.

This is a concrete example of liquid surface value being exchanged for illiquid, long-duration capacity.

Macro Scaling

Under the locked model, the same exchange operates at system scale once \(\lambda(t)\) approaches its terminal floor. The residual stock (anchored to the $3.65 T factual baseline) is forced into funding that exchange, expanding the long-duration claim layer that later feeds the Concentration Trap and the G-SIB cascade.

Status

Reciprocal mirror = accounting conservation rule for physical monetary value exchange.  
No thermodynamic phrasing remains.  
Corporate CapEx/FCF compression serves as the observable micro analogue.  
The 2055/2059 timeline and $25–45 T liquidation range remain outcomes of the residual baseline + locked dynamics under zero intervention.

The wording is now aligned with standard financial and accounting notation.
